\SourcePage{259}{262}
\section{Wave functions of particles with arbitrary spin}

Having developed the formal algebra of spinors of arbitrary rank, we can now
turn to the immediate problem: the properties of wave functions for particles
of arbitrary spin.

It is convenient to approach this question by considering a collection of
$n$ particles of spin $1/2$. The largest possible value of the $z$ component
of the system's total spin is $n/2$. It is obtained when $s_z=1/2$ for every
particle (all spins point in the same direction, along the $z$ axis). In this
case the total spin $S$ of the system must likewise equal $n/2$.

\SourcePage{260}{263}
Every component of the particle-system wave function
$\psi(\sigma_1,\sigma_2,\ldots,\sigma_n)$ then vanishes except for
$\psi(1/2,1/2,\ldots,1/2)$. Write the wave function as a product of $n$
spinors $\psi^\lambda\varphi^\mu\ldots$, one for each particle. Each spinor
has only its component with $\lambda,\mu,\ldots=1$ different from zero.
Consequently, the sole nonzero product is $\psi^1\varphi^1\ldots$. The full
set of these products, however, constitutes a spinor of rank $n$ symmetric in
all its indices. A transformation of coordinates (so that the spins no longer
point along the $z$ axis) yields a general spinor of rank $n$, which remains
symmetric.

The spin properties of wave functions are, in essence, their behavior under
rotations of the coordinates. They are therefore identical for a particle of
spin $s$ and for a system of $n=2s$ spin-$1/2$ particles oriented so that the
system has total spin $s$. It follows that the wave function of a spin-$s$
particle is a symmetric spinor of rank $n=2s$.

The number of independent components of a symmetric spinor of rank $2s$ is,
as required, $2s+1$. Indeed, the distinct components are precisely those
whose indices contain $2s$ ones and no twos, then $2s-1$ ones and one two,
and so on, through no ones and $2s$ twos.

Mathematically, symmetric spinors classify the possible transformation types
of quantities under rotations of the coordinates. Suppose that $2s+1$
different quantities transform linearly among themselves, and that no choice
of linear combinations can reduce their number. Their transformation law is
then equivalent to that of the components of a symmetric spinor of rank $2s$.
Any collection of functions, of whatever size, that transform linearly among
themselves under coordinate rotations can be reduced, by a suitable linear
transformation, to one or more symmetric spinors\footnote{Equivalently,
symmetric spinors realize the irreducible representations of the rotation
group (see \S~98).}.

Thus an arbitrary spinor of rank $n$, $\psi^{\lambda\mu\nu\ldots}$, can be
reduced to symmetric spinors of ranks $n$, $n-2$, $n-4$, \ldots. In practice
the reduction may be performed as follows. Symmetrizing
$\psi^{\lambda\mu\nu\ldots}$ over all its indices produces a symmetric
spinor of the same rank $n$. Next, contracting

\SourcePage{261}{264}
the original spinor $\psi^{\lambda\mu\nu\ldots}$ over various pairs of
indices gives spinors of rank $n-2$ such as
$\psi^\lambda{}_{\lambda}{}^{\nu\ldots}$. Symmetrizing these in turn gives
symmetric spinors of rank $n-2$. Contracting
$\psi^{\lambda\mu\ldots}{}_{\lambda\mu}$ over two pairs and symmetrizing the
resulting spinors gives symmetric spinors of rank $n-4$, and so forth.

It remains to connect the components of a symmetric spinor of rank $2s$ with
the $2s+1$ functions $\psi(\sigma)$, where
$\sigma=s,s-1,\ldots,-s$. The component
\[
 \psi^{\overbrace{11\ldots1}^{s+\sigma}
              \overbrace{22\ldots2}^{s-\sigma}},
\]
in which the index 1 occurs $s+\sigma$ times and the index 2 occurs
$s-\sigma$ times, corresponds to a spin projection $\sigma$ on the $z$ axis.
To verify this, replace the one particle of spin $s$ once again by a system of
$n=2s$ spin-$1/2$ particles. The component just written corresponds to the
product
\[
 \underbrace{\psi^1\varphi^1\ldots}_{s+\sigma}
 \underbrace{\chi^2\rho^2\ldots}_{s-\sigma};
\]
this product represents a state in which $s+\sigma$ particles have spin
projection $+1/2$ and $s-\sigma$ have projection $-1/2$. Their total
projection is therefore
$\frac12(s+\sigma)-\frac12(s-\sigma)=\sigma$. Finally, choose the
proportionality factor between this spinor component and $\psi(\sigma)$ so
that
\begin{equation}
 \sum_{\sigma=-s}^{+s}|\psi(\sigma)|^2
 =\sum_{\lambda,\mu,\ldots=1}^{2}|\psi^{\lambda\mu\ldots}|^2
 \tag{57.1}
\end{equation}
(This sum is a scalar, as it must be, since it gives the probability of finding
the particle at the specified point in space.) In the sum on the right-hand
side, a component having $(s+\sigma)$ indices equal to 1 occurs
\[
 \frac{(2s)!}{(s+\sigma)!(s-\sigma)!}
\]
times. It follows that the correspondence between the functions $\psi(\sigma)$
and the spinor components is given by
\begin{equation}
 \psi(\sigma)=\left[\frac{(2s)!}{(s+\sigma)!(s-\sigma)!}\right]^{1/2}
 \psi^{\overbrace{11\ldots1}^{s+\sigma}
       \overbrace{22\ldots2}^{s-\sigma}}.
 \tag{57.2}
\end{equation}


\SourcePage{262}{265}
Relation (57.2) ensures not only condition (57.1), but also, as is readily
checked, the more general condition
\begin{equation}
 \psi^{\lambda\mu\ldots}\varphi_{\lambda\mu\ldots}
 =\sum_\sigma(-1)^{s-\sigma}\psi(\sigma)\varphi(-\sigma),
 \tag{57.3}
\end{equation}
where $\psi^{\lambda\mu\ldots}$ and $\varphi_{\lambda\mu\ldots}$ are two
different spinors of equal rank, while $\psi(\sigma)$ and $\varphi(\sigma)$
are the functions associated with them by (57.2). The factor
$(-1)^{s-\sigma}$ appears because raising all the indices changes the sign of
a component once for every index equal to 2.

Equations (55.5) specify the action of the spin operator on the wave functions
$\psi(\sigma)$. We can readily determine how these operators act when the wave
function is written as a spinor of rank $2s$. For spin $1/2$, the functions
$\psi(+1/2)$ and $\psi(-1/2)$ coincide with the spinor components $\psi^1$ and
$\psi^2$. According to (55.6) and (55.7), the action of the spin operators on
them gives
\begin{equation}
 \begin{aligned}
 (\hat s_x\psi)^1&=\tfrac12\psi^2,&
 (\hat s_y\psi)^1&=-\tfrac i2\psi^2,&
 (\hat s_z\psi)^1&=\tfrac12\psi^1,\\
 (\hat s_x\psi)^2&=\tfrac12\psi^1,&
 (\hat s_y\psi)^2&=\tfrac i2\psi^1,&
 (\hat s_z\psi)^2&=-\tfrac12\psi^2.
 \end{aligned}
 \tag{57.4}
\end{equation}

To extend the construction to an arbitrary spin, consider once more a system
of $2s$ particles of spin $1/2$, and express its wave function as a product of
$2s$ spinors. The spin operator for this system is the sum of the individual
particles' spin operators. Each term acts only on its corresponding spinor,
with the result specified by formulas (57.4). Returning then to arbitrary
symmetric spinors, that is, to the wave functions of a particle of spin $s$,
gives
\begin{equation}
 \begin{aligned}
 (\hat s_x\psi)^{\overbrace{11\ldots}^{s+\sigma}
                         \overbrace{22\ldots}^{s-\sigma}}
 &=\frac{s+\sigma}{2}\psi^{\overbrace{11\ldots}^{s+\sigma-1}
                                  \overbrace{22\ldots}^{s-\sigma+1}}
 +\frac{s-\sigma}{2}\psi^{\overbrace{11\ldots}^{s+\sigma+1}
                                  \overbrace{22\ldots}^{s-\sigma-1}},\\
 (\hat s_y\psi)^{\overbrace{11\ldots}^{s+\sigma}
                         \overbrace{22\ldots}^{s-\sigma}}
 &=-i\frac{s+\sigma}{2}\psi^{\overbrace{11\ldots}^{s+\sigma-1}
                                   \overbrace{22\ldots}^{s-\sigma+1}}
 +i\frac{s-\sigma}{2}\psi^{\overbrace{11\ldots}^{s+\sigma+1}
                                  \overbrace{22\ldots}^{s-\sigma-1}},\\
 (\hat s_z\psi)^{\overbrace{11\ldots}^{s+\sigma}
                         \overbrace{22\ldots}^{s-\sigma}}
 &=\sigma\psi^{\overbrace{11\ldots}^{s+\sigma}
                       \overbrace{22\ldots}^{s-\sigma}}.
 \end{aligned}
 \tag{57.5}
\end{equation}

Until now, spinors have been discussed as wave functions for the intrinsic
angular momentum of elementary particles.  Formally, however, the spin of a
single particle is no different from the total angular momentum of any system
regarded as a whole, when its internal structure is set aside.

\SourcePage{263}{266}
Consequently, the transformation properties of spinors apply equally to the
behaviour under spatial rotations of the wave functions $\psi_{jm}$ for any
particle (or particle system) having total angular momentum $j$, irrespective
of whether that angular momentum is orbital or intrinsic.  There must
therefore be a definite correspondence between the transformation laws of the
eigenfunctions $\psi_{jm}$ under rotations of the coordinate system and those
of the components of a symmetric spinor of rank $2j$.


In establishing this correspondence, two aspects of the dependence of the
wave functions on the angular-momentum projection $m$, for fixed $j$, must be
distinguished carefully. A wave function may be regarded as a probability
amplitude for the possible values of $m$, or it may be an eigenfunction for a
specified value of $m$.

Both aspects appeared at the beginning of \S~55 in the eigenfunction
$\delta_{\sigma\sigma_0}$ of $s_z$ belonging to the value
$s_z=\sigma_0$. Their mathematical difference is particularly transparent
for a particle with spin $s=1/2$. With respect to the variable $\sigma$, its
spin function is a contravariant spinor of rank one and should therefore be
written, in spinor notation, as $\delta^\lambda{}_{\sigma_0}$. With respect
to $\sigma_0$, it is consequently a covariant spinor.

This observation is evidently quite general: formulas analogous to (57.2)
establish a correspondence between the eigenfunctions $\psi_{jm}$ and the
components of a covariant symmetric spinor of rank $2j$\footnote{The same
result can also be reached by a somewhat different argument.  Expand the
wave function $\psi$ of a particle whose angular momentum is $j$ in the
eigenfunctions $\psi_{jm}$:
$\psi=\sum_m a_m\psi_{jm}$.  The coefficients $a_m$ are then the probability
amplitudes associated with the different values of $m$.  In this sense, they
correspond to the ``components'' $\psi(m)$ of the spin wave function, and
this correspondence fixes their transformation law.  On the other hand, the
value of $\psi$ at a specified point in space cannot depend on the coordinate
system chosen; hence the sum $\sum a_m\psi_{jm}$ must be a scalar.  Comparison
with the scalar in (57.3) shows that the $a_m$ must transform in the same way
as $(-1)^{j-m}\psi_{j,-m}$.}:
\begin{equation}
 \psi_{jm}=\left[\frac{(2j)!}{(j+m)!(j-m)!}\right]^{1/2}
 \psi_{\overbrace{11\ldots}^{j+m}\overbrace{22\ldots}^{j-m}}.
 \tag{57.6}
\end{equation}
For integral angular momentum $j$, the eigenfunctions are the spherical
harmonics $Y_{jm}$.  The case $j=1$ is particularly important.


\SourcePage{264}{267}
The three spherical functions $Y_{1m}$ are
\[
 Y_{10}=i\sqrt{\frac3{4\pi}}\cos\theta=i\sqrt{\frac3{4\pi}}n_z,
 \qquad
 Y_{1,\pm1}=\mp i\sqrt{\frac3{8\pi}}\sin\theta e^{\pm i\varphi}
 =\mp i\sqrt{\frac3{8\pi}}(n_x\pm in_y)
\]
where $\mathbf n$ is the unit vector in the radius-vector direction. Their
transformation properties show that these three functions are equivalent to
the components of a vector $\mathbf a$ under the correspondence
\begin{equation}
 \psi_{10}=ia_z,\qquad
 \psi_{11}=-\frac i{\sqrt2}(a_x+ia_y),\qquad
 \psi_{1,-1}=\frac i{\sqrt2}(a_x-ia_y).
 \tag{57.7}
\end{equation}
Comparing these expressions with (57.6) gives the following correspondence
between the components of a symmetric spinor of rank two and those of a
vector:
\begin{equation}
 \psi_{11}=-\frac i{\sqrt2}(a_x+ia_y),\qquad
 \psi_{12}=\psi_{21}=\frac i{\sqrt2}a_z,\qquad
 \psi_{22}=\frac i{\sqrt2}(a_x-ia_y),
 \tag{57.8}
\end{equation}
or
\begin{equation}
 \psi^{11}=\frac i{\sqrt2}(a_x-ia_y),\qquad
 \psi^{12}=\psi^{21}=-\frac i{\sqrt2}a_z,\qquad
 \psi^{22}=-\frac i{\sqrt2}(a_x+ia_y).
 \tag{57.9}
\end{equation}
Conversely,
\begin{equation}
 a_z=i\sqrt2\psi^{12},\qquad
 a_x=\frac i{\sqrt2}(\psi^{22}-\psi^{11}),\qquad
 a_y=\frac1{\sqrt2}(\psi^{11}+\psi^{22}).
 \tag{57.10}
\end{equation}
With these definitions one readily verifies
\begin{equation}
 \psi^{\lambda\mu}\varphi_{\lambda\mu}=\mathbf a\cdot\mathbf b,
 \tag{57.11}
\end{equation}
where $\mathbf a$ and $\mathbf b$ are the vectors corresponding to the
symmetric spinors $\psi^{\lambda\mu}$ and $\varphi_{\lambda\mu}$. One may
also verify the spinor--vector correspondence\footnote{The mixed components
of a symmetric spinor may be written as $\psi^\lambda{}_{\mu}$ without
distinguishing $\psi^{\lambda\mu}$ from $\psi^{\mu\lambda}$.}
\begin{equation}
 \psi^\lambda{}_{\rho}\varphi^{\rho\mu}
 -\varphi^\lambda{}_{\rho}\psi^{\rho\mu}
 \longleftrightarrow i\sqrt2(\mathbf a\times\mathbf b).
 \tag{57.12}
\end{equation}
Equations (57.10) can be written compactly in terms of the Pauli matrices:
\begin{equation}
 a_i=\frac i{\sqrt2}(\hat\sigma_i)^\lambda{}_{\mu}\psi^\mu{}_{\lambda}.
 \tag{57.13}
\end{equation}

\SourcePage{265}{268}
The matrix indices on $\sigma$ have been placed above and below in accordance
with the positions of the spinor indices on $\psi^\mu{}_{\lambda}$. The origin
of this formula becomes clear in the special case in which the rank-two
spinor $\psi^\mu{}_{\lambda}$ reduces to the product of a rank-one spinor
$\psi^\mu$ and its complex conjugate $\psi_\lambda^*$. The quantity
\[
 \frac12\psi_\lambda^*\hat{\boldsymbol\sigma}^\lambda{}_{\mu}\psi^\mu
\]
is then the mean spin for a particle with wave function $\psi^\mu$, so its
vector character is evident.

The correspondence (57.8) or (57.9) is a special case of a general rule. To
every symmetric spinor of even rank $2j$, with integral $j$, one may associate
a symmetric tensor of half that rank, $j$, whose contraction over every pair
of indices vanishes; such a tensor will be called irreducible. This follows
already from the equality of the numbers of independent components of these
spinors and tensors, both being $2j+1$, as a direct count confirms\footnote{In
other words, the $2j+1$ components of an irreducible tensor of rank $j$ (for
integral $j$), the set of $2j+1$ spherical functions $Y_{jm}$, and the
$2j+1$ components of a symmetric spinor of rank $2j$ all realize the same
irreducible representation of the rotation group.}. The relation between the
spinor and tensor components can be found from (57.8)--(57.10), by treating a
spinor of the given rank as a product of rank-two spinors and the tensor as a
product of vectors.

\ProblemHeading

1. Rewrite the definition (57.4) of the spin-$1/2$ operator in terms of the
spinor components of the vector $\mathbf s$.

\SolutionHeading

Using (57.9), which relates the vector $\mathbf s$ to the spinor
$s^{\lambda\mu}$, definition (57.4) becomes
\[
 s^{\lambda\mu}\psi^\nu=\frac i{2\sqrt2}
 (\psi^\lambda g^{\mu\nu}+\psi^\mu g^{\lambda\nu}).
\]

2. Give the formulas defining the action of the spin operator on the vector
wave function of a particle of spin 1.

\SolutionHeading

Equations (57.9) relate the components of the vector function
$\boldsymbol\psi$ to the components of the spinor $\psi^{\lambda\mu}$, while
the last of equations (57.5) gives
\[
 \hat s_z\psi_+=-\psi_+,
 \qquad \hat s_z\psi_-=\psi_-,
 \qquad \hat s_z\psi_z=0
\]
where $\psi_\pm=\psi_x\pm i\psi_y$, or
\[
 \hat s_z\psi_x=-i\psi_y,\qquad
 \hat s_z\psi_y=i\psi_x,\qquad
 \hat s_z\psi_z=0.
\]
The remaining formulas follow by cyclic permutation of the indices $x,y,z$.
Together they read
\[
 \hat s_i\psi_k=-ie_{ik\ell}\psi_\ell.
\]

The complex vector $\boldsymbol\psi$ may be represented as
$\boldsymbol\psi=e^{i\alpha}(\mathbf u+i\mathbf v)$, where $\mathbf u$ and
$\mathbf v$ are real vectors. By a suitable choice of the common phase
$\alpha$, these vectors may be made mutually perpendicular. The two vectors
$\mathbf u$ and $\mathbf v$ specify a plane with the following property: the
spin projection on the direction normal to it can take only the values
$\pm1$.
